\section{Joukowski descent and the arithmetic of reciprocal levels}
\label{sec:joukowski-descent}

The reciprocal equation admits a useful quotient which separates its
complex-analytic geometry from its arithmetic level.  Put
\[
 J(z)=z+z^{-1}\qquad(z\in\C^\times).
\]
The involution $z\mapsto z^{-1}$ is the deck involution of $J$.  In this
section we descend the reciprocal $\Ein$-sum through $J$, determine the
critical points of the descended entire function, and then apply
Theorem~\ref{thm:master} to its values at algebraic points.

\subsection{The entire quotient and its critical points}

Let $C_0(w)=2$, $C_1(w)=w$, and
\[
 C_{n+1}(w)=wC_n(w)-C_{n-1}(w)\qquad(n\geq1).
\]
Thus $C_n(w)=2T_n(w/2)$ and
\begin{equation}\label{eq:joukowski-chebyshev}
 C_n(J(z))=z^n+z^{-n}.
\end{equation}

\begin{proposition}[Entire Joukowski descent]\label{prop:entire-joukowski-descent}
There is a unique entire function $H$ such that
\begin{equation}\label{eq:H-descent}
 H(J(z))=\Ein(z)+\Ein(z^{-1})\qquad(z\in\C^\times).
\end{equation}
It is given by the locally uniformly convergent series
\begin{equation}\label{eq:H-chebyshev-series}
 H(w)=\sum_{n\geq1}\frac{(-1)^{n-1}}{n\,n!}C_n(w).
\end{equation}
Moreover
\begin{align}
 H'(w)
 &=2e^{-w/2}
   \frac{\sinh\!\left(\frac12\sqrt{w^2-4}\right)}{\sqrt{w^2-4}}
   \label{eq:Hprime-sinh}\\
 &=e^{-w/2}\prod_{n\geq1}
   \left(1+\frac{w^2-4}{4\pi^2n^2}\right).
   \label{eq:Hprime-product}
\end{align}
The right side of \eqref{eq:Hprime-sinh} is understood by its removable
continuation at $w=\pm2$; in particular $H'(2)=e^{-1}$ and
$H'(-2)=e$.
\end{proposition}

\begin{proof}
If $K\subset\C$ is compact, there is an $R_K>1$ such that both roots of
$X^2-wX+1$ have modulus at most $R_K$ for every $w\in K$.  Hence
\eqref{eq:joukowski-chebyshev} gives
$|C_n(w)|\leq2R_K^n$ on $K$.  The series
\eqref{eq:H-chebyshev-series} therefore converges locally uniformly and
defines an entire function.  Substitution of $w=J(z)$ and the power series
of $\Ein$ proves \eqref{eq:H-descent}.  Uniqueness follows because $J$ is
surjective: the equation $z^2-wz+1=0$ has a nonzero root for every
$w\in\C$.

For $z\ne\pm1$, differentiation of \eqref{eq:H-descent} gives
\[
 H'(J(z))
 =\frac{e^{-1/z}-e^{-z}}{z-z^{-1}}.
\]
Write $w=z+z^{-1}$ and $v=z-z^{-1}$, so $v^2=w^2-4$.  Then
\[
 e^{-1/z}-e^{-z}
 =e^{-w/2}\bigl(e^{v/2}-e^{-v/2}\bigr),
\]
which proves \eqref{eq:Hprime-sinh} away from $w=\pm2$.  The quotient
$2\sinh(v/2)/v$ is an entire even function of $v$, so the apparent
singularities are removable.  Finally Euler's product
\[
 \frac{\sinh u}{u}=\prod_{n\geq1}
   \left(1+\frac{u^2}{\pi^2n^2}\right)
\]
with $u=\frac12\sqrt{w^2-4}$ gives
\eqref{eq:Hprime-product}.
\end{proof}

\begin{corollary}[Critical lattice and the real diffeomorphism]
\label{cor:H-critical-lattice}
The critical points of $H$ are exactly the simple points
\begin{equation}\label{eq:H-critical-points}
 \xi_n^\pm=\pm2i\sqrt{\pi^2n^2-1}\qquad(n\geq1).
\end{equation}
In particular every critical point is transcendental.  On the real axis
$H'(x)>0$, and $H$ maps $\mathbb R$ increasingly and bijectively onto
$\mathbb R$.
\end{corollary}

\begin{proof}
The product \eqref{eq:Hprime-product} gives
\eqref{eq:H-critical-points}; each zero is simple because it occurs in
exactly one factor and is nonzero.  If one of these points were algebraic,
then $\pi^2=((\xi_n^\pm)^2/(-4)+1)/n^2$ would be algebraic, a contradiction.
For real $x$, every factor in \eqref{eq:Hprime-product} is positive, since
\[
 1+\frac{x^2-4}{4\pi^2n^2}
 \geq1-\frac1{\pi^2n^2}>0.
\]
Thus $H'(x)>0$.

For $w\to+\infty$, choose the root $z>1$ of $J(z)=w$.  Then
$z=w+O(w^{-1})$ and
\[
 H(w)=\Ein(z)+\Ein(z^{-1})
     =\gamma+\log w+O(w^{-1})\longrightarrow+\infty.
\]
For $w\to-\infty$, write the root of modulus greater than one as $z=-a$,
$a\to+\infty$.  Since
$\Ein(-a)=\gamma+\log a-\Ei(a)$ and
$\Ei(a)\sim e^a/a$, one has $H(w)\to-\infty$.  Strict monotonicity now
proves the last assertion.
\end{proof}

\begin{proposition}[A rational differential equation]
\label{prop:H-ode}
The function $H$ satisfies
\begin{equation}\label{eq:H-third-order-ode}
 w(w^2-4)H'''
 +(w^3+2w^2-4w+4)H''
 +(w^2-w+2)H'=0.
\end{equation}
Consequently every additive translate $H-2c$ satisfies the same equation.
\end{proposition}

\begin{proof}
Put $P(w)=e^{w/2}H'(w)$.  By \eqref{eq:Hprime-sinh},
\[
 P(w)=2\frac{\sinh(\frac12\sqrt{w^2-4})}{\sqrt{w^2-4}}.
\]
Direct differentiation gives
\[
 w(w^2-4)P''+2(w^2+2)P'-\frac{w^3}{4}P=0.
\]
Substitution of $P=e^{w/2}H'$ and collection of terms gives
\eqref{eq:H-third-order-ode}.  The equation contains no undifferentiated
$H$, so it is invariant under additive translation.
\end{proof}

\subsection{Algebraic values after descent}

For $w\in\Qbar$, let
\[
 \Lambda(w)=\{\lambda\in\Qbar^\times:
                    \lambda+\lambda^{-1}=w\}
\]
as a set.  At $w=2$ or $w=-2$ this set has one element, but the identity
\eqref{eq:H-descent} counts that element twice.

\begin{theorem}[Multipoint arithmetic descent]
\label{thm:H-multipoint-descent}
Let $w_1,\ldots,w_s\in\Qbar$ be distinct, let
$\Lambda=\bigcup_{j=1}^s\Lambda(w_j)$, and put
\[
 d=\dim_{\Q}\Span_{\Q}\Lambda.
\]
Then
\begin{equation}\label{eq:H-multipoint-trdeg}
 \operatorname{trdeg}_{\Qbar}
 \Qbar\bigl(\{e^\lambda\}_{\lambda\in\Lambda},
             H(w_1),\ldots,H(w_s)\bigr)=d+s.
\end{equation}
In particular $H(w_1),\ldots,H(w_s)$ are algebraically independent over
the field generated by the exponential coordinates.  Hence $H(w)$ is
transcendental for every algebraic $w$.
\end{theorem}

\begin{proof}
By Theorem~\ref{thm:master}, the coordinates
$Y_\lambda=\Ein(\lambda)$, $\lambda\in\Lambda$, are algebraically
independent over
$E=\Qbar(\{e^\lambda\}_{\lambda\in\Lambda})$, and
$\operatorname{trdeg}_{\Qbar}E=d$.  Each $H(w_j)$ is a nonzero linear form
in the $Y_\lambda$ supported on $\Lambda(w_j)$, with coefficient $2$ at a
double root.  The supports for distinct $j$ are disjoint.  These $s$ linear
forms can therefore be extended to a linear coordinate system in the
$Y_\lambda$, and are algebraically independent over $E$.  This proves
\eqref{eq:H-multipoint-trdeg}.  The last assertion is the case $s=1$.
\end{proof}

\begin{corollary}[Algebraic levels have transcendental preimages]
\label{cor:H-algebraic-levels}
For every $c\in\Qbar$, every solution of
\begin{equation}\label{eq:H-algebraic-level}
 H(w)=2c
\end{equation}
is transcendental.  Equivalently, every solution of
\begin{equation}\label{eq:reciprocal-algebraic-level}
 \Ein(z)+\Ein(z^{-1})=2c
\end{equation}
is transcendental.
\end{corollary}

\begin{proof}
An algebraic solution of \eqref{eq:H-algebraic-level} would contradict the
last assertion of Theorem~\ref{thm:H-multipoint-descent}.  If $z$ in
\eqref{eq:reciprocal-algebraic-level} were algebraic, then
$w=J(z)$ would be algebraic and would solve \eqref{eq:H-algebraic-level}.
\end{proof}

\begin{corollary}[The unique real orbit of a real level]
\label{cor:H-real-levels}
For every $c\in\mathbb R$ there is a unique $w_c\in\mathbb R$ with
$H(w_c)=2c$.  It satisfies
\[
 \begin{array}{c|c|c}
 \text{condition on }c&\text{location of }w_c&
 \text{preimage under }J\\ \hline
 \Ein(-1)<c<\Ein(1)&-2<w_c<2&
 \text{one conjugate pair on }|z|=1,\\
 c>\Ein(1)&w_c>2&\text{two positive reciprocal points},\\
 c<\Ein(-1)&w_c<-2&\text{two negative reciprocal points}.
 \end{array}
\]
At $c=\Ein(1)$ or $c=\Ein(-1)$, respectively, $z=1$ or $z=-1$ is a
double zero of \eqref{eq:reciprocal-algebraic-level}.  If $c$ is algebraic,
the points in the table are all transcendental.

More generally, for arbitrary real $c$, every nonreal algebraic solution
of $H(w)=2c$ is impossible.  Thus at most the unique real preimage $w_c$
can be algebraic.
\end{corollary}

\begin{proof}
Existence and uniqueness follow from Corollary~\ref{cor:H-critical-lattice}.
The thresholds follow from
$H(2)=2\Ein(1)$ and $H(-2)=2\Ein(-1)$, and the three descriptions are the
elementary inverse images of the intervals $(-2,2)$, $(2,\infty)$ and
$(-\infty,-2)$ under $J$.  Since
\[
 J(z)-2=\frac{(z-1)^2}{z},\qquad
 J(z)+2=\frac{(z+1)^2}{z},
\]
and $H'(\pm2)\ne0$, the endpoint zeros have order two.  The assertion for
algebraic $c$ is Corollary~\ref{cor:H-algebraic-levels}.

Finally, if a nonreal algebraic $w$ solved $H(w)=2c$ with $c$ real, then
$\bar w$ would be a distinct algebraic point with
$H(\bar w)=\overline{H(w)}=2c$.  This contradicts the algebraic
independence of $H(w)$ and $H(\bar w)$ in
Theorem~\ref{thm:H-multipoint-descent}.
\end{proof}

\subsection{The Euler level and witness amplification}

Put
\[
 G_\gamma(w)=H(w)-2\gamma.
\]

\begin{theorem}[Euler level in the quotient]
\label{thm:H-euler-level}
The function $G_\gamma$ has exactly one real zero $w_*$, and
\begin{equation}\label{eq:w-star-location}
 -2<w_*<2,\qquad w_*=2\cos\theta_*,
\end{equation}
A high-precision Newton iteration gives the non-rigorous decimal values
\[
 w_*=1.0073046182687490652\ldots,
 \qquad \theta_*=1.0429750684400255695\ldots.
\]
The two preimages of $w_*$ are the unit-circle zeros
$\alpha_*=e^{i\theta_*}$ and $\overline{\alpha_*}$ of the reciprocal
$\Ein$ equation.  Every algebraic zero of $G_\gamma$ must equal $w_*$;
equivalently every algebraic zero of the reciprocal equation belongs to
$\{\alpha_*,\overline{\alpha_*}\}$.

If $w_*$, equivalently $\alpha_*$, is algebraic, then $\gamma$ and
$\delta$ are algebraically independent.  More precisely, put
$a=\alpha_*$, $b=a^{-1}$,
\[
 D_*=\Ein(a)-\Ein(b),\qquad
 d_*=\dim_{\Q}\Span_{\Q}\{1,a,b\}.
\]
Then
\begin{equation}\label{eq:witness-amplification}
 \operatorname{trdeg}_{\Qbar}
 \Qbar(e,e^a,e^b,\gamma,\delta,D_*)=d_*+3,
\end{equation}
and $\gamma,\delta,D_*$ are algebraically independent over
$\Qbar(e,e^a,e^b)$.
\end{theorem}

\begin{proof}
Since
\[
 H(-2)=2\gamma-2\Ei(1)<2\gamma
 <2\gamma+2E_1(1)=H(2),
\]
Corollary~\ref{cor:H-critical-lattice} gives the unique zero $w_*$ and
\eqref{eq:w-star-location}.  If an algebraic zero $w$ were nonreal, the
last assertion of Corollary~\ref{cor:H-real-levels} would exclude it; if it
were real, uniqueness would force $w=w_*$.  Passing between $w$ and the
roots of $z^2-wz+1$ proves the equivalent statement in the $z$-plane.

Now assume that $w_*$ is algebraic.  Then $a$ and $b$ are distinct
algebraic numbers.  Apply Theorem~\ref{thm:master} to
$\Lambda=\{1,a,b\}$.  Over
$E=\Qbar(e,e^a,e^b)$ the three numbers
\[
 Y_1=\Ein(1),\qquad Y_a=\Ein(a),\qquad Y_b=\Ein(b)
\]
are algebraically independent.  The reciprocal equation and
$\delta=eE_1(1)$ give
\[
 \gamma=\frac{Y_a+Y_b}{2},\qquad
 \frac{\delta}{e}=Y_1-\frac{Y_a+Y_b}{2},\qquad
 D_*=Y_a-Y_b.
\]
This is an invertible linear change of the three $Y$-coordinates; replacing
$\delta/e$ by $\delta$ is invertible over $E$.  This proves
\eqref{eq:witness-amplification}.
\end{proof}

\begin{corollary}[Conditional completion of the six-constant field]
\label{cor:witness-six-constant-field}
If $\alpha_*$ is algebraic, then
\[
 \operatorname{trdeg}_{\Qbar}
 \Qbar\bigl(\gamma,\delta,E_1(1),\Ei(1),\Ci(1),\Chi(1)\bigr)=5.
\]
The complete ideal of algebraic relations among these six displayed
coordinates is generated by
\begin{equation}\label{eq:six-constant-only-relation}
 \Ei(1)-E_1(1)-2\Chi(1)=0.
\end{equation}
\end{corollary}

\begin{proof}
Use the notation
\[
 A=\Ein(1),\qquad B=\Ein(-1),\qquad
 C=\frac{\Ein(i)+\Ein(-i)}2.
\]
The point $\alpha_*$ is distinct from $\pm1,\pm i$.  Indeed
$w_*\in(-2,2)$ excludes $\pm1$, while
\[
 H(0)=2\gamma-2\Ci(1)<2\gamma
\]
excludes $\pm i$; here
$\Ci(1)=\gamma+\int_0^1(\cos t-1)\,dt/t>\gamma-1/4>0$.
Theorem~\ref{thm:master}, applied to
$\{1,-1,i,-i,\alpha_*,\alpha_*^{-1}\}$, now shows that
\[
 \gamma,\quad e,\quad A,\quad B,\quad C
\]
are algebraically independent: they are one exponential coordinate and
four independent linear forms with disjoint supports in the relevant
$\Ein$-coordinates.  The six-constant field equals
$\Qbar(\gamma,e,A,B,C)$, so its transcendence degree is five.

Identity \eqref{eq:six-constant-only-relation} follows from
$E_1(1)=A-\gamma$, $\Ei(1)=\gamma-B$, and
$\Chi(1)=\gamma-(A+B)/2$.  The kernel of the evaluation map from a
polynomial ring in the six displayed coordinates is a height-one prime.
It contains the irreducible linear polynomial in
\eqref{eq:six-constant-only-relation}; hence it is precisely the principal
ideal generated by that polynomial.
\end{proof}

\subsection{Every level has asymptotic transcendental fibres}

The following strengthens the existence part of Picard's theorem for this
specific function: there is no exceptional finite level, and an explicit
sequence of preimages is available at every level.

\begin{theorem}[Asymptotic level fibres]
\label{thm:H-level-asymptotics}
Fix $c\in\C$.  For every sufficiently large integer $k$ there is a simple
solution $w_{k,c}$ of $H(w)=2c$ such that
\begin{equation}\label{eq:H-level-zero-asymptotic}
 w_{k,c}
 =-\log(2\pi k)-\log\log(2\pi k)
  +i\left(2\pi k+\frac\pi2\right)
  +O_c\!\left(\frac1{\log k}\right).
\end{equation}
Consequently every finite value of $H$ is assumed infinitely often.  If
$c$ is algebraic, every member of these fibres is transcendental.  If
$c\in\mathbb R$, the conjugates $\overline{w_{k,c}}$ give a second family.
\end{theorem}

\begin{proof}
It is convenient first to work in the $z$-plane.  Put
\[
 \mathcal F_c(z)=\Ein(z)+\Ein(z^{-1})-2c
                =H(J(z))-2c.
\]
On $\C\setminus(-\infty,0]$, with principal logarithm, set
\[
 L_c(z)=\Log z+\gamma-2c.
\]
The connection formula for $\Ein$ gives
\begin{equation}\label{eq:Fc-connection}
 \mathcal F_c(z)=L_c(z)+E_1(z)+\Ein(z^{-1}).
\end{equation}
Uniformly in a fixed closed sector about the positive imaginary axis,
the standard expansion of $E_1$ \cite{DLMF} and the Taylor expansion of
$\Ein$ give
\begin{equation}\label{eq:Fc-asymptotic-input}
 E_1(z)=\frac{e^{-z}}z\bigl(1+O(z^{-1})\bigr),
 \qquad
 \Ein(z^{-1})=z^{-1}+O(z^{-2}).
\end{equation}

Let
\[
 R_k=2\pi k,\qquad
 a_k=\log R_k+\log\log R_k,\qquad
 \zeta_k=-a_k+i\left(R_k+\frac\pi2\right).
\]
On $D_k=\{z:|z-\zeta_k|\leq1\}$, for all sufficiently large $k$, the
principal branches of $\Log z$ and $\Log L_c(z)$ are analytic.  Uniformly
on $D_k$,
\begin{align}
 \Log z&=\log R_k+\frac{\pi i}{2}
          +O(a_k/R_k),\label{eq:log-z-on-Dk}\\
 \Log L_c(z)&=\log\log R_k+O_c((\log R_k)^{-1}).
 \label{eq:log-L-on-Dk}
\end{align}
Define
\[
 \Phi_{k,c}(z)=(2k+1)\pi i-\Log z-\Log L_c(z).
\]
Equations \eqref{eq:log-z-on-Dk}--\eqref{eq:log-L-on-Dk} give
\[
 \Phi_{k,c}(z)=\zeta_k+O_c((\log R_k)^{-1}),
 \qquad
 \Phi_{k,c}'(z)=-\frac1z-\frac1{zL_c(z)}=O(R_k^{-1}).
\]
Thus $\Phi_{k,c}$ maps $D_k$ into itself and is a contraction for large
$k$.  Its fixed point $u_{k,c}$ satisfies
\begin{equation}\label{eq:model-root-asymptotic}
 u_{k,c}=\zeta_k+O_c((\log R_k)^{-1})
\end{equation}
and
\[
 u_{k,c}e^{u_{k,c}}L_c(u_{k,c})=-1.
\]
Hence $u_{k,c}$ is a zero of the model function
\[
 h_c(z)=1+ze^zL_c(z).
\]

Write $q_c(z)=z+\Log z+\Log L_c(z)$, so
$h_c=1+e^{q_c}$ on $D_k$.  At $u_{k,c}$,
\begin{equation}\label{eq:model-derivative}
 q_c'(u_{k,c})
 =1+\frac1{u_{k,c}}+\frac1{u_{k,c}L_c(u_{k,c})}
 =1+O(R_k^{-1}).
\end{equation}
In particular the model zero is simple.

Now multiply the exact function by its nonvanishing normalizing factor:
\[
 A_c(z)=ze^z\mathcal F_c(z).
\]
From \eqref{eq:Fc-connection}--\eqref{eq:Fc-asymptotic-input},
\begin{equation}\label{eq:exact-model-error}
 A_c(z)-h_c(z)
 =ze^z\left(E_1(z)-\frac{e^{-z}}z\right)
  +ze^z\Ein(z^{-1}).
\end{equation}
The first term is $O(R_k^{-1})$ near $u_{k,c}$.  The model equation gives
\[
 |e^{u_{k,c}}|
 =\frac1{|u_{k,c}|\,|L_c(u_{k,c})|}
 =O_c((R_k\log R_k)^{-1}),
\]
so the second term in \eqref{eq:exact-model-error} is
$O_c((R_k\log R_k)^{-1})$.  Consequently there is a constant $C_c$ such
that
\begin{equation}\label{eq:exact-model-uniform-bound}
 |A_c(z)-h_c(z)|\leq \frac{C_c}{R_k}
\end{equation}
uniformly on $|z-u_{k,c}|\leq1$ for all sufficiently large $k$.

Fix $M>2C_c$.  Equation \eqref{eq:model-derivative} and Taylor's theorem
give, on $|z-u_{k,c}|=M/R_k$,
\[
 |h_c(z)|\geq\frac{M}{2R_k}
\]
for large $k$.  A second application of the same Taylor estimate, now
comparing $h_c$ with
$-q_c'(u_{k,c})(z-u_{k,c})$, shows that $h_c$ has exactly one zero, counted
with multiplicity, in this circle.  By
\eqref{eq:exact-model-uniform-bound} and Rouch\'e's theorem, $A_c$ also has
exactly one zero there.  Since $ze^z$ does not vanish, this is a simple zero
$z_{k,c}$ of $\mathcal F_c$, and
\[
 z_{k,c}=u_{k,c}+O_c(R_k^{-1}).
\]

Finally put $w_{k,c}=J(z_{k,c})$.  Since
$z_{k,c}^{-1}=O(R_k^{-1})$ and $J'(z_{k,c})=1+O(R_k^{-2})$, this is a
simple zero of $H(w)-2c$ and has the same asymptotic
\eqref{eq:H-level-zero-asymptotic}.  The arithmetic assertion follows from
Corollary~\ref{cor:H-algebraic-levels}, and conjugation gives the lower
family for real $c$.
\end{proof}

\begin{corollary}[Order and a zero-count lower bound]
\label{cor:H-order-zero-count}
For every $c\in\C$, the entire function $H-2c$ has order one and
exponential type one.  If $n_c(r)$ counts its zeros in $|w|\leq r$, with
multiplicity, then
\[
 \frac{r}{\pi}-O_c(\log r)\leq n_c(r)=O_c(r).
\]
\end{corollary}

\begin{proof}
Formula \eqref{eq:Hprime-sinh} gives exponential type at most one, and the
negative-real asymptotic used in Corollary~\ref{cor:H-critical-lattice}
gives type at least one.  Jensen's formula then gives $n_c(r)=O_c(r)$.
For real $c$, Theorem~\ref{thm:H-level-asymptotics} and conjugation give the
lower bound directly.  For general $c$, apply the same contraction argument
with negative indices as well as positive ones; equivalently repeat the
proof in a fixed sector about the negative imaginary axis.  The two
families have modulus $2\pi k+O_c(\log k)$, which gives the stated bound.
\end{proof}

\subsection{The derivative-level de Branges structure}

Define
\begin{equation}\label{eq:K-rotated-derivative}
 K(t)=e^{it/2}H'(it).
\end{equation}

\begin{proposition}[Rotated derivative]
\label{prop:H-rotated-LP}
One has
\begin{align}
 K(t)
 &=\frac{2\sin\!\left(\frac12\sqrt{t^2+4}\right)}{\sqrt{t^2+4}}
 \label{eq:K-sine}\\
 &=\sin(1)\prod_{n\geq1}
   \left(1-\frac{t^2}{4(\pi^2n^2-1)}\right).
 \label{eq:K-real-product}
\end{align}
Thus $K$ belongs to the Laguerre--P\'olya class and has only the simple real
zeros
\[
 \pm2\sqrt{\pi^2n^2-1}\qquad(n\geq1).
\]
Furthermore
\begin{equation}\label{eq:HB-function}
 E(t)=K(t)+iK'(t)
\end{equation}
is a Hermite--Biehler function and therefore defines a de Branges space.
\end{proposition}

\begin{proof}
Equations \eqref{eq:K-sine} and \eqref{eq:K-real-product} follow from
\eqref{eq:Hprime-sinh} and Euler's product for $\sin u/u$; the constant is
$\prod_{n\geq1}(1-1/(\pi^2n^2))=\sin(1)$.  The partial products in
\eqref{eq:K-real-product} are real-rooted polynomials and converge locally
uniformly, which is exactly the required Laguerre--P\'olya approximation.

For $\Im t>0$, the logarithmic derivative of the product satisfies
\[
 \Im\frac{K'(t)}{K(t)}
 =\sum_{n\geq1}\Im\left(
   \frac1{t-a_n}+\frac1{t+a_n}\right)<0,
 \qquad a_n=2\sqrt{\pi^2n^2-1}.
\]
The paired series converges locally uniformly off the real zero set.  With
$E^\#(t)=\overline{E(\bar t)}=K(t)-iK'(t)$, it follows that
\[
 |E(t)|^2-|E^\#(t)|^2
 =-4|K(t)|^2\Im\frac{K'(t)}{K(t)}>0
 \qquad(\Im t>0).
\]
This is the Hermite--Biehler inequality.
\end{proof}

\begin{remark}[Scope of the real-zero analogy]
The terminology above is classical in the sense of P\'olya--Schur and de
Branges.\footnote{G.~P\'olya and J.~Schur, \emph{Über zwei Arten von
Faktorenfolgen in der Theorie der algebraischen Gleichungen}, J. Reine
Angew. Math. 144 (1914), 89--113; L.~de Branges, \emph{Hilbert Spaces of
Entire Functions}, Prentice--Hall, 1968.}
It applies one derivative below the reciprocal-level problem.  Every
translate $H-2c$ has the same derivative, hence the same $K$, the same
Hermite--Biehler function \eqref{eq:HB-function}, and the same canonical de
Branges space.  This construction therefore forgets the level $c$ and in
particular cannot distinguish $c=\gamma$ from another real level.

Conversely, Theorem~\ref{thm:H-level-asymptotics} gives infinitely many
nonreal zeros of $H-2c$.  For real $c$, the function $H-2c$ is therefore not
in the Laguerre--P\'olya class and cannot be the real part of a
Hermite--Biehler function, whose real and imaginary parts have only real
zeros.  Likewise, although the symmetry $z\leftrightarrow z^{-1}$ resembles
a self-inversive or Lee--Yang symmetry, the off-circle fibres show that no
full circle theorem holds here.  The statement that algebraic zeros at the
Euler level must lie on the unit circle is an arithmetic restriction, not a
Lee--Yang theorem.
\end{remark}

\begin{remark}[Novelty and logical status]
The analytic descent and product formulas above are elementary consequences
of the Joukowski symmetry and Euler's sine product.  The arithmetic
statements are corollaries of Theorem~\ref{thm:master}; their novelty is the
packaging, not a new transcendence mechanism.  The conditional hypothesis
that $\alpha_*$ is algebraic is not known, so
Theorem~\ref{thm:H-euler-level} does not prove the transcendence of
$\gamma$ or $\delta$.  No claim of priority should be made for the
Joukowski packaging before a specialist literature search.
\end{remark}
